\subsection{Example 16-3a: FRM Exam 1999 - Question 83 (Fat tails)}

\textbf{Enunciado}

\emph{In the presence of fat tails in the distribution of returns, VAR based on
the delta-normal method would (for a linear portfolio)}
\begin{enumerate}[label=\alph*.]
\item \hl{Underestimate the true VAR}
\item Be the same as the true VAR
\item Overestimate the true VAR
\item Cannot be determined from the information provided
\end{enumerate}

\subsubsection*{Solución}


\textbf{Qué hace el delta-normal. }Para un portafolio lineal, el método delta-normal supone P\&L aproximadamente normal y calcula
\[
\VaR_\alpha^{DN}\approx z_\alpha \,\sigma,
\]
donde \(\sigma\) es la desviación estándar del P\&L y \(z_\alpha\) es el cuantil normal.

\textbf{Efecto de colas pesadas. }``Fat tails'' implica más probabilidad de pérdidas extremas que bajo normalidad, manteniendo varianza comparable.
Por ello, para niveles altos \(\alpha\), el cuantil verdadero suele ser más grande que el cuantil normal con la misma \(\sigma\):
\[
\VaR_\alpha^{true} > \VaR_\alpha^{DN}.
\]


El delta-normal \textbf{subestima} el VaR verdadero:
(a) Underestimate the true VAR.










